% ============================================================
%  Chapter 3 — Wi-Fi and Routers
%  The Secret Life of the Internet
% ============================================================

\chapter{Wi-Fi and Routers}

\chapteropening{%
  How does your tablet get to the internet if there is no wire plugged in?
  Willa Wi-Fi and Rory Router have the answer!
}
\chapterbadge{Waves, Signals, and Router Magic}

% --- Opening story ---
\section*{The Invisible Bridge}

Before Dot could zoom off through the internet, there was a problem.

Sam's tablet was sitting on the sofa.
There was no wire connecting it to anything.

\begin{center}
  \Large\bfseries\color{WarmOrange}
  ``How does the signal get OUT of my tablet?''
\end{center}

A friendly shimmer appeared in the air.
It was Willa Wi-Fi — not a person exactly, more like a warm glow of colourful waves.

\character{Willa Wi-Fi}{That's me!  I carry signals through the air, like invisible ripples on a pond.
  You can't see me, but I'm always here, helping your device reach the router.}

\vspace{6pt}

\begin{picturethis}
  Drop a pebble in a still pond.
  Ripples spread out in rings from the centre.
  Wi-Fi works just like that —
  only instead of water ripples, it uses \textbf{radio waves} spreading from the router.
  Your tablet can feel those waves and use them to send and receive information.
\end{picturethis}

\sectionrule

% --- What is Wi-Fi ---
\section*{What Is Wi-Fi?}

\term{Wi-Fi} is a technology that lets devices connect to a network using radio waves instead of cables.
The radio waves carry information back and forth between your device and a \term{router}.

Wi-Fi is very convenient — you can move around a room, wander to the kitchen, or sit in the garden,
and as long as you are close enough to the router, you stay connected.

Common devices that use Wi-Fi:
\begin{itemize}
  \item Tablets and smartphones
  \item Laptops
  \item Smart TVs
  \item Game consoles
  \item Smart speakers
\end{itemize}

\begin{funfact}
Wi-Fi uses radio waves on two main frequencies — 2.4 GHz (travels further) and 5 GHz (faster but shorter range). Your router might use both at the same time!
\end{funfact}

\sectionrule

% --- What is a router? ---
\section*{Meet Rory Router}

Down the hall, sitting on a shelf, was a small box with little blinking lights.
That was Rory Router.

\character{Rory Router}{Hi there!  I'm the one who directs traffic.
  Every time a device in this house wants to reach the internet,
  the message comes to me first.
  I figure out the best path and send it on its way.}

A \term{router} is a device that sits in your home (or school, or cafe) and does two key jobs:

\begin{enumerate}[label=\textbf{\arabic*.}]
  \item It connects all the devices in your home together into a \term{local network}.
  \item It connects that local network to the wider internet through a cable that leads to your \term{ISP}.
\end{enumerate}

\begin{picturethis}
  Imagine Rory Router as the postmaster of a small village.
  Every letter (piece of information) that comes into the village
  passes through Rory.
  Rory reads the address and sends each letter to exactly the right home.
  And when villagers want to send letters to the outside world,
  Rory posts them on their behalf.
\end{picturethis}

\sectionrule

% --- Journey step by step ---
\section*{The Journey From Tablet to Internet}

\begin{quickmap}
  \textbf{How Dot leaves Sam's tablet and reaches the internet:}
  \begin{enumerate}[label=\textbf{\arabic*.}]
    \item Sam presses play on a video.
    \item The tablet sends a request as radio waves through Wi-Fi.
    \item Willa Wi-Fi carries those waves to the router.
    \item Rory Router receives the request and checks where it needs to go.
    \item Rory sends the request out through a cable to the ISP.
    \item The ISP forwards it to the wider internet.
    \item The answer (the video data) travels back the same way.
    \item The tablet receives the data and the video plays!
  \end{enumerate}
\end{quickmap}

\begin{diagram}[A home network with one router connecting many devices.]
  \resizebox{\linewidth}{!}{%
  \begin{tikzpicture}[>=Stealth]
    \node[draw=LeafGreen!80!black, rounded corners, fill=LeafGreen!20, minimum width=2.4cm, minimum height=1cm] (router) at (0,0) {Router};
    \node[draw=InternetBlue, rounded corners, fill=SkyBlue!20, minimum width=2.3cm, minimum height=0.85cm] (tablet) at (0,2.4) {Tablet};
    \node[draw=WarmOrange!80!black, rounded corners, fill=WarmOrange!20, minimum width=2.3cm, minimum height=0.85cm] (laptop) at (2.7,1.2) {Laptop};
    \node[draw=CoralRed!80!black, rounded corners, fill=CoralRed!16, minimum width=2.3cm, minimum height=0.85cm] (phone) at (2.7,-1.2) {Phone};
    \node[draw=SoftPurple!80!black, rounded corners, fill=SoftPurple!18, minimum width=2.3cm, minimum height=0.85cm] (tv) at (0,-2.4) {Smart TV};
    \node[draw=SunYellow!70!black, rounded corners, fill=SunYellow!25, minimum width=2.8cm, minimum height=0.85cm] (speaker) at (-2.8,0) {Smart Speaker};
    \node[draw=InternetBlue!80!black, rounded corners, fill=InternetBlue!12, minimum width=2.5cm, minimum height=1cm] (cloud) at (6.1,0) {Internet};

    \draw[dashed,thick,InternetBlue] (router) -- node[midway,above,sloped]{Wi-Fi} (tablet);
    \draw[dashed,thick,InternetBlue] (router) -- node[midway,above,sloped]{Wi-Fi} (laptop);
    \draw[dashed,thick,InternetBlue] (router) -- node[midway,below,sloped]{Wi-Fi} (phone);
    \draw[dashed,thick,InternetBlue] (router) -- node[midway,below,sloped]{Wi-Fi} (tv);
    \draw[dashed,thick,InternetBlue] (router) -- node[midway,above,sloped]{Wi-Fi} (speaker);
    \draw[thick,->,InternetBlue] (router) -- (cloud);
  \end{tikzpicture}%
  }
\end{diagram}

\sectionrule

% --- Wi-Fi vs wired ---
\section*{Wi-Fi vs a Cable — What's the Difference?}

Some devices connect to the router using a physical cable called an \term{Ethernet cable}.
This is often faster and more reliable than Wi-Fi, but you have to stay in one place.

\begin{center}
  \begin{tabularx}{\linewidth}{>{\bfseries}l Y Y}
    Feature    & Wi-Fi                & Ethernet cable \\
    Movement   & Move around freely   & Must stay near the cable \\
    Speed      & Very fast            & Usually faster \\
    Reliability & Can drop in some spots & Very stable \\
    Common use & Phones, tablets, laptops & Desktop computers, smart TVs \\
  \end{tabularx}
\end{center}

\sectionrule

\begin{newwords}
  \begin{description}[labelwidth=3.2cm, leftmargin=3.4cm]
    \item[\term{Wi-Fi}]         Wireless technology that uses radio waves to connect devices to a network.
    \item[\term{Router}]        A device that connects your home network to the internet and directs traffic.
    \item[\term{Radio waves}]   Invisible energy waves that carry Wi-Fi signals through the air.
    \item[\term{Local network}] The private network inside your home that connects your devices together.
    \item[\term{Ethernet}]      A cable that connects a device directly to a router.
  \end{description}
\end{newwords}

\sectionrule

\begin{tryit}
  \textbf{Trace the path!}

  On paper, draw:
  \begin{itemize}
    \item A tablet in the bedroom.
    \item A router in the hallway.
    \item Wavy lines from the tablet to the router (those are Wi-Fi radio waves!).
    \item A solid line from the router to a big cloud labelled \textbf{The Internet}.
  \end{itemize}

  Now label each part:
  \emph{device, Wi-Fi signal, router, ISP cable, internet}.

  Can you add a phone, a laptop, and a TV all connecting to the same router?
\end{tryit}

\sectionrule

\begin{bigidea}
  \textbf{Wi-Fi} carries signals through the air using radio waves.
  The \textbf{router} receives those signals and connects your home
  to the wider internet — directing traffic in and out.
\end{bigidea}

\character{Dot}{We've left the house!
  Now we're zooming toward a huge building full of powerful computers.
  Let's go meet Sunny Server!}
